% -------------------------------------
%               Setup
% -------------------------------------
\documentclass[t, 24pt]{beamer}
\usepackage{lipsum}
\usepackage{listings} % for code
\usepackage{etoolbox} % For \ifstrempty00
\newcommand{\weekcount}{3} % \def\weekcount{3}

% Setup > KU Setup
\usepackage{KUstyle}
\toplinje{Meeting | Week \weekcount} % The text at top. Remove the command if no text is desired

% Code
\usepackage{fancyvrb}
\VerbatimEnvironment

% -------------------------------------
%               Code
% -------------------------------------
\definecolor{bgcolor}{RGB}{248,248,248} % Light gray background
\definecolor{keywordcolor}{RGB}{0,0,160} % Dark blue keywords
\definecolor{stringcolor}{RGB}{163,21,21} % Dark red strings
\definecolor{commentcolor}{RGB}{0,128,0} % Dark green comments
\definecolor{numbercolor}{RGB}{128,0,128} % Purple line numbers
\definecolor{promptcolor}{RGB}{64,64,64} % Dark gray prompts

% Custom style for Python code
\lstdefinestyle{professionalPython}{
    language=Python,
    backgroundcolor=\color{bgcolor},
    basicstyle=\linespread{1.1}\tiny\ttfamily,
    keywordstyle=\color{keywordcolor}\bfseries,
    commentstyle=\color{commentcolor}\itshape,
    stringstyle=\color{stringcolor},
    numberstyle=\tiny\color{numbercolor},
    identifierstyle=\color{black},
    tabsize=4,
    showstringspaces=false,
    numbers=left,
    stepnumber=1,
    numbersep=10pt,
    frame=single,
    framerule=0.5pt,
    rulecolor=\color{gray},
    breaklines=true,
    captionpos=b,
    xleftmargin=15pt,
    xrightmargin=5pt,
    aboveskip=10pt,
    belowskip=-10pt,
    morekeywords={def, return, class, from, import, as, if, elif, else, while, for, in, try, except, finally, with, yield, lambda, global, nonlocal, assert, pass, break, continue, raise},
    literate=*{:}{{\textcolor{keywordcolor}{:}}}1,
}


% -------------------------------------
% Command with question slides counter
% -------------------------------------

\newcounter{subjectNumber}
\setcounter{subjectNumber}{0}

\makeatletter
\newcommand{\resetsubject}{%
\stepcounter{subjectNumber}%
\@ifundefined{c@currentQSlide\thesubjectNumber}{%
\newcounter{currentQSlide\thesubjectNumber}%
}{}%
\setcounter{currentQSlide\thesubjectNumber}{0}%
}

\newcommand{\incrementSlideCnt}{%
\stepcounter{currentQSlide\thesubjectNumber}%
}

\AtEndDocument{%
\@tempcnta=1 \loop\ifnum\@tempcnta<\numexpr\value{subjectNumber}+1 \immediate\write\@auxout{%
\string\@ifundefined{c@totalQSlide\the\@tempcnta}{%
\string\newcounter{totalQSlide\the\@tempcnta}%
}{}%
\string\setcounter{totalQSlide\the\@tempcnta}{\the\value{currentQSlide\the\@tempcnta}}%
}%
\advance\@tempcnta by1 \repeat }
\makeatother



\newenvironment{numberedSlide}[2][]{%
  % Choose frame option
  \ifstrempty{#1}{%
    \begin{frame}[hoved]
  }{%
    \begin{frame}[#1,hoved]
  }%
  \incrementSlideCnt
  % Fallback check for totalQSlide
  \frametitle{%
    #2 \space
    \ifcsname c@totalQSlide\thesubjectNumber\endcsname (\arabic{currentQSlide\thesubjectNumber}/\arabic{totalQSlide\thesubjectNumber})
    \else
      (\arabic{currentQSlide\thesubjectNumber}/?)
    \fi
  }
}{%
  \end{frame}
}



%\newenvironment{numberedSlide}[1]{%
%  \begin{frame}[hoved]
%    \incrementSlideCnt
%    \frametitle{#1 (\arabic{currentQSlide\thesubjectNumber}/\arabic{totalQSlide\thesubjectNumber})}
%}{%
%  \end{frame}
%}


% -------------------------------------
%                Setup
% -------------------------------------
\resetsubject % Initialize the first subject

% -------------------------------------
%                TiKz
% -------------------------------------

\usepackage{tikz}
\usetikzlibrary{shapes.geometric, arrows}

\tikzstyle{startstop}
=
[rectangle,
rounded
corners,
minimum
width=3.5cm,
minimum
height=1cm,
text
centered,
draw=black,
fill=blue!20]
\tikzstyle{process}
=
[rectangle,
minimum
width=3.5cm,
minimum
height=1cm,
text
centered,
draw=black,
fill=green!20]
\tikzstyle{arrow}
=
[thick,->,>=stealth]

\begin{document}
    % Intro slides.
    { \setbeamertemplate{background}{\includegraphics[width=\paperwidth,height=\paperheight]{KU/forside.pdf}} \begin{frame}\begin{textblock*}{\textwidth}(0\textwidth,0.1\textheight) \begin{beamercolorbox}[wd=6.3cm,ht=7.7cm,sep=0.5cm]{hvidbox} \fontsize{4}{10}\fontfamily{ptm}\selectfont \textls[200]{UNIVERSITY OF COPENHAGEN} \noindent\textcolor{KUrod}{\rule{5.3cm}{0.4pt}}\end{beamercolorbox}\end{textblock*} \begin{textblock*}{\textwidth}(0\textwidth,0.1\textheight) \begin{beamercolorbox}[wd=6.3cm,sep=0.5cm]{hvidbox} \Large \textcolor{KUrod}{Meeting - Week \weekcount} \vspace{0.5cm} \par \large Progress, Challenges \& Next steps \vspace{0.5cm} \par \textcolor{gray}{\scriptsize Simon Winther \& Hjalte Bjoernstrup, Department of Computer Science, University of Copenhagen (DIKU), \today.}\end{beamercolorbox}\end{textblock*} \begin{textblock}{1}(6,11.44) \includegraphics[width=1cm]{KU/KU-logo.png}\end{textblock}\end{frame} }

    % ----------------------------------------------------
    %                       Slide 1
    % ----------------------------------------------------
    \begin{numberedSlide}{Questions?} 
        Should we train a separate model for each organ (e.g., liver, brain, prostate) to maximize specialization, or would a single model trained on all datasets generalize well across different organs? Alternatively, would pretraining on multiple organs and fine-tuning per organ be a better approach?
        \\~\\
        The LiTS preprocessing script resizes images to 320×320 and creates three-channel
        inputs using neighboring slices.
        \begin{itemize}
            \item Should we crop/resize our dataset to match this format, or can
                we keep the original resolution?

            \item How should we handle three-channel input if the number of slices
                is not a multiple of 3?

            \item Are there any strict formatting requirements we need to follow
                for U-Net 3+?
        \end{itemize}
    \end{numberedSlide}

    % ----------------------------------------------------
    %                       Slide 2
    % ----------------------------------------------------
    \begin{numberedSlide}{Questions?} 
        \textbf{Handling 4D NIfTI Files in U-Net 3+}

        Some .nii files contain an extra dimension, making them 4D (e.g.,
        512×512×49×4).

        \begin{itemize}
            \item How should we handle the extra dimension? Should we select a
                specific channel, merge them, or treat them as separate inputs?

            \item If selecting a single channel, which one is most relevant for
                segmentation?

            \item Should we preprocess 4D files differently from standard 3D volumes?
        \end{itemize}

        \textbf{How do we fairly compare this U-net 3+, with our multi-scaling
        approach if we need to crop images differently/more in U-net 3+, than
        our multi-scaling model approach, since then the different cropping seem
        to hold a big factor for the final score.}
    \end{numberedSlide}

    % ----------------------------------------------------
    %                       Slide 3
    % ----------------------------------------------------
\begin{numberedSlide}{Questions?}
\begin{itemize}
\item It seems like we have test data with images and labels and then validation data without labels. Why is that? Is it, because we need to get the individual score of each of the segmentation and average it out as the average dice score of all images?  

\item We're unsure if we have to calculate the dice coefficient on the decoded segmentation and the mask? Since it doesn't seem like we have a mask in the validation part?

\end{itemize}
    
\end{numberedSlide}


\end{document}
